
1900年到1956年
数学家铺平理论道路

20世纪初的四位数学家用方程推动整个世界。

David Hilbert 及希尔伯特23问题
哥德尔、图灵冯诺依曼  维纳



希尔伯特在20世纪之初的演讲为几十年内的数学突飞猛进提供了指路牌。

20世纪之初的希尔伯特23问题，为数学后来的突飞猛进及人工智能的萌芽提供了指路牌。

1928年国际数学家大会上，希尔伯特提出一个宏伟计划：运用公理化方法统一整个数学。
希望能严谨证明任意公理系统内的所有命题彼此相容无矛盾，即在形式化的算术系统内部：

1. 完备性：一切数学命题都可通过公理集有限步推定真伪
2. 相容性：不能从公理系统导出矛盾，如不可能得到1+1=2且1+1=3
3. 可判定性 (希尔伯特第十问题)： 给定任一丢番图方程，设计一个过程/方法，通过有限次的计算，判定该方程是否可解.



希尔伯特纲领：

①**公理化**数学。

②**完备性：**所有数学内容都能够在这个公理系统中被证真或证伪。

③**相容性：**可证明这个公理系统是相容的/无矛盾的。

④**可判定性：**存在一个算法来判定是否一个语句能够被证明。



希尔伯特期望数学的各个分支都能建立完全的公设集，从而将数学建立在完美严谨的逻辑基石上。“利用这种新的数学基础，我将可以解决世界上所有的基础性问题”。
计划有一定的进展，几乎全世界的数学家都乐观地看着数学大厦即将竣工。

逻辑美梦只做了三年，打击接二连三。



希尔伯特被陆续打脸：

p1931年，哥德尔不完备定理“包含了算术的数学整体不可能既完备又相容”：

Ø哥德尔第一定理(否定完备性)：任何包含自然数的相容形式系统F中，存在不可判定命题S，即S在F中既不能被证真，也不能被证伪

Ø哥德尔第二定理(否定相容性)：任何包含自然数系的相容形式系统 F，其相容性/无矛盾性不能在 F中被证明。

p1936年，图灵à存在数学问题不能用任何机械过程来解决

Ø可计算性=一组明确、可以执行步骤的有序操作，在有限的时间内终止，并产生结果。

Ø可计算性 = 图灵机能解决

Ø图灵机的停机问题不可判定

Ø图灵机不能解决的问题>>>>图灵机能解决的问题





1931年，哥德尔不完备性定理“包含了算术的数学整体不可能既完备又相容”：

1. 哥德尔第一定理(否定完备性)：对于包含自然数系的任何相容的形式系统F，存在F中的不可判定命题S，在这个系统中既不能被证实，也不能被证伪。（数论中迄今为止既没有证明也没有推翻的各种猜想是不可判定的吗？）
2. 哥德尔第二定理(否定相容性)：对于包含自然数系的任何相容的形式系统 F，F的相容性/无矛盾性不能在 F中被证明。

1936年，Church,Turing 分别否定希尔伯特判定问题

1. 美国普林斯顿大学的Alonzo Church教授：基于自创的λ演算
2. 英国剑桥国王学院读本科的图灵：图灵机的停机问题不可判定



哥德尔（Kurt Godel，1906-1978）出生于捷克斯洛伐克。
1924年进入维也纳大学，主修物理和数学，开始学习逻辑。1940年1月到达美国普林斯顿。
1948年入籍仪式中，哥德尔发现美国宪法存在逻辑漏洞，有向独裁专制演变的逻辑可能性，试图与法官辩论，被作为入籍见证人的爱因斯坦和摩根坦极力制止。
哥德尔是20世纪最伟大的数学家和逻辑学家。爱因斯坦把哥德尔在数学中的贡献与他本人对物理学的贡献相提并论。
美国《时代》杂志曾经评选出对20世纪思想产生重大影响的100人中，哥德尔被列为第四位。



**哥德尔不完备性定理的划时代意义-数学**

**哥德尔不完备性定理使数学、逻辑学发生革命性的变化：**

- Ø摧毁了数学的所有重要领域能被完全公理化这一强烈的信念
- Ø扑灭了沿着希尔伯特曾设想的路线证明数学内部相容性的全部希望
- Ø使得人们不得不必须重新评价普遍认可的数学哲学。
- 

现代科学在哲学方面的三大成果：

​    哥德尔不完备性定理

​	塔尔斯基的形式语言的真理论

​	图灵机判定问题



**哥德尔不完备性定理的划时代意义**-科学**

哥德尔不完备性定理的科学和哲学价值远超出数学领域，启发后人对哲学本质、世界基本问题的思考：
不完备性定理表明在同一个系统中完备性和逻辑自洽不可兼得：
一个理论中的完备性与一致性不相容，才提供了理论体系进一步发展的突破口
一个理论要求的完备性，可能存在于下一个更深层的理论中。
例如，欧几里得几何最后被“非欧几何”所完备；
           牛顿力学和经典电磁论最后被量子力学和相对论在更深的层次 “完备”。

一致性与完备性不可兼得 v.s.量子物理中海森堡不确定性原理“动量、位置不能同时确定”：不完备性与不确定性在哲学上勾画出了人类知识的疆界和认识的极限。这两个原理之于人类的意义超过了牛顿力学、万有引力、相对论等，后者可能影响几个世纪，而前者会影响整个人类的文明历史。
爱因斯坦的相对论颠覆传统时空观、哥德尔不完备性定理令逻辑基础轰然倒塌、测不准原理引发新一轮“真实是否可被感觉”的哲学大讨论，被称为20世纪理论大跃进的三个支柱。



** **四大数学家之图灵**

阿兰·图灵（Alan Turing)，英国著名数学家和逻辑学家，被称为计算机科学之父、人工智能之父，是计算机逻辑的奠基者，提出了“图灵机”和“图灵测试”等重要概念。
曾协助英国军方破解德国的著名密码系统Enigma(哑谜机)，帮助盟军取得了二战的胜利。
“图灵奖”的设立是为了纪念他在计算机领域的卓越贡献。



希尔伯特第10问题--判定丢番图方程的可解性：“给定一个系数均为有理整数，包含任意个未知数的丢番图方程，设计一个过程(process)去，通过有限次的计算，能够判定该方程在整数上是否可解。” 
（是否有明确程序在有限时间内可判定丢番图方程是否可解？）
 “过程”即指“算法”—通过有限多的步骤，对数学函数进行有效计算的方法
什么样的数学函数可以有效计算—数理逻辑学对可计算理论的发展
数学归纳法
数理逻辑中研究递归函数的递归论/可计算性理论



**图灵机****--****计算机的理论原型**



图灵的成就在计算机领域是里程碑式的：
 严格定义了“明确程序”的概念
 提出的图灵机为电子计算机的发明奠定了基础
 证明了计算存在局限



**计算的本质**

**算法：**一组明确、可以执行步骤的有序操作，在有限的时间内终止，并产生结果。

20世纪30年代以前，人们并没有真正认识计算的本质
古代中国的“演算法”：一个数学问题，只有可用算盘解算时，才算可解。
自图灵机思想提出不到10年，世界上第一台电子计算机诞生了。
现代计算机反映计算的抽象、理论、设计过程：
抽象：将问题抽象/形式化为可计算问题
理论：寻找解决可计算问题的方法
设计：模型的具体实现问题

**问题1：求数组逆序数**

**问题2：数组排序问题**



**算法理论**

特性：有效性、确定性、可终止性、
分类：递归算法、迭代算法、穷举算法、贪心算法、……
评价：时间复杂度、空间复杂度
P问题与NP(Non-deterministic Polynomial)问题是非确定性多项式问题
自动机理论
分布式算法、并行算法



从计算角度认知思维、视觉和生命过程



符号主义者à认知是一种符号处理过程，因此认知就是计算

视觉认知理论者à视觉（ 听觉、嗅觉、……）是一种计算

DNA计算技术的可行性à 生命过程也是一种计算

那么，现在的各种LLM可以说具有生命吗？



**认知世界****-****我们对自身了解多少？**

<!-- ![截屏2024-07-07 17.31.51](/Users/Min369/Library/Application Support/typora-user-images/截屏2024-07-07 17.31.51.png) -->





**四大数学家之冯·诺依曼**

美籍匈牙利数学家、计算机科学家、物理学家。 
一位现代科学全才：
语言能力：精通匈牙利语、英文、德语、法语等7门语言
记忆力：能逐字逐句把读过一本书或一篇文章默写下来，还能把《双城记》立即从原来的语言直接翻译成英语
存储量：博览群书，自由爱好历史，读遍了绝大多数百科全书式的历史书
计算：6岁心算两个八位数的除法、8岁掌握微积分、能在头脑中编写和修改计算算法



有趣的是，这些能力如今轻易被计算机实现。

AI已经成为我们的衍生外脑。天才的标准如今不再适用！



**天赋v.s.成就**



一种观点：与“天才中的天才”相比，他的成就，认真算下来，只能算是“碌碌无为”，其实是不匹配的。
数学量子力学流体力学博弈论参与原子弹研制推动计算机的发明
没有攻克世纪难题，开创数学分支，拿菲尔兹奖
没有深入研究量子力学，取得开创性贡献，获得诺奖
现代计算机之父”、“博弈论之父” 
数学家外尔：“以前冯·诺依曼数学搞得多么好，如今却如此不务正业。”
爱因斯坦：“我不能容忍这样的科学家，他拿出一块木板来，寻找最薄的地方，然后在容易钻透的地方钻许多孔。”
快与慢的抉择



**电子计算机的发明—二战的巨大压力**



在英国，图灵 •Colossus：
动机：提高破译德军密电的速度
方案：把机械的“炸弹”机改成更快速的电路装置1943年，Colossus(巨人)研制成功。
效果：将破译密码的时间由几周缩短到几小时。
功绩：助攻诺曼底成功登陆---二战转折点。



在美国，冯·诺依曼 •ENIAC：
动机：研制原子弹涉及超几十亿次的数学运算和逻辑指令，人力难以完成。
方案：将大型电子数字积分机改进为能作各种科学计算的通用计算机。1946年，ENIAC建造成功。
效果：将原本几个月的人力计算缩短到几天。
功绩：为世界上第一颗氢弹的研制立下汗马功劳。



共同特点：利用打孔卡输入、运用真空管计算、
		   体积庞大、对二战胜利功勋卓著
相似缺陷：只为专门目的设计，不能储存程序。
改进方向：能完成任何目的的图灵通用机。



**后续的故事：二战后**



在英国：
战时无节制的战争经费战后经济危机和官僚作风
强大的协作需求 v.s. 孤僻的性格
图灵：造计算机的难点主要是硬件而非数学模型，就把琐碎的工程问题留给工程师吧。
“想出来”就是“做出来”
图灵开始思考机器是否能够具备智能。



在美国：
未受战争破坏，ENIAC掀起的计算机和电子工程科学大发展
二战期间，希特勒给罗斯福送上的“大礼”---爱因斯坦、费米、冯·诺依曼、冯·卡门、哥德尔、塔斯基、外尔、……等顶级数学家、物理学家
在科技上长于应用而弱于基础大翻身
美苏冷战“核威慑“一扇关不上的大门” 



**美国制霸电子信息技术**

50年代，冯·诺依曼帮助IBM的小沃森完成了第一套存储程序计算机的开发。
1957年冯· 诺伊曼逝世后，普林斯顿高等研究院退出了计算机科学后续发展。



全球顶尖人才开始源源不断涌入：近80年的美国科技高速发展

<!-- ![截屏2024-07-07 17.34.57](/Users/Min369/Library/Application Support/typora-user-images/截屏2024-07-07 17.34.57.png) -->



**机器智能**

“依据什么标准评价一台机器是否具备智能？ ”--- 1950年，图灵《机器能思考吗？》，提出 “图灵测试”。
他开始为虚拟计算机“想”下象棋的程序。
他扮演虚构计算机，严格执行自己的程序，进行象棋比赛。一步耗时半小时。



**四大数学家之维纳****(****Norbert Wiener)**

诺伯特·维纳，著名应用数学家，控制论创始人。
维纳在70年的科学生涯中，先后涉足哲学、数学、物理、生物学，在各领域中都取得了丰硕成果。
他先后留学于英国剑桥大学和德国哥丁根大学，在罗素、哈代、希尔伯特等著名数学家指导下研究逻辑和数学。



1948年，发表《控制论——关于在动物和机器中控制和通讯的科学》提出了控制论：
认为机器与人的统一性—人和机器都是通过反馈来实现某种目的，揭示了用机器模拟人的可能性，为人工智能的提出奠定了重要基础。



1950年，维纳出版《人有人的用处：控制论与社会》，将自动机器和人进行了对比，认为两者均由
	感知装置	
	信息传递装置
	行动装置
  构成一个相对独立的系统。
这一系统接受、处理、储存和传递信息，并做出一系列行动，以实现与系统之外世界的互动。



**图灵和维纳的人工智能思想**

图灵认为实现人工智能需要大致三步：
Step1[智能的表述] 使用一种通用语言将现实任务描述成特定的数学命题。
Step2[智能的设计] 找到一种算法解决这种通用语言描述的数学命题问题。
Step3[智能的测试] 找到一种算法验证所提算法的正确性。
解决Step 1的“描述性问题”并不比Step 2的“解决问题”更加容易。
图灵的想法接近功能主义，着眼于智能机器内部，主要研究智能机器的内在结构和运行方式，把机器和环境的关系放在次要地位。
维纳的想法接近行为主义，强调智能机器和环境的关系，着眼于外部，研究机器与环境各种输入输出的关系以及根据这些关系来调节机器。



图灵和维纳的思想都对人工智能的研究之路产生了重要影响。
图灵方法和维纳方法的区别：
图灵设想的人工智能是一次性成型的
维纳设想的人工智能是尝试-测试-调整：自主寻找，反复迭代多次后成型。
维纳的思想为在更少人类智能参与的情况下，实现人工智能带来了新的解决途径。

<!-- ![image-20240707173643550](/Users/Min369/Library/Application Support/typora-user-images/image-20240707173643550.png) -->



**维纳思想：神经网络**

p践行维纳思想获得成功的研究：神经网络的反向传播算法

p**(****深度****)****神经网络获得的巨大成功表明，维纳的****尝试****—****测试****—****调整框架****具有很强的灵活性和适用性**。

p最近20年来，更值得关注的维纳思想成功应用是强化学习：

①定义机器的可能状态和所能执行行动的集合；

②借鉴行为主义理论，界定了行动的回报函数，设定任务完成所获得的回报最大。

③机器按照某种简单规则来探索不同行动获得的回报，验证当前机器所采取的行动策略是否需要优化

④更新机器下一步所采取的行动策略

⑤反复迭代直至任务完成。



p强化学习首先在维纳创立的控制理论中取得了丰硕成果，接着在人工智能、机器学习的很多领域中都获得了成功应用。

p近年来，强化学习和深度学习相结合，使用深度神经网络来储存所习得知识，籍此推测未验证的情况应采取何种行动。而这一点又和维纳曾经非常关注的神经网络和联结主义研究联系了起来。

pAlphaGo的成功表明，采用尝试—测试—调整框架可以解决一些之前认为难以解决的难题。
